\section{M\'ascaras de bits}
\whenuse{Representar subconjuntos peque\~nos, enumerar subm\'ascaras o guardar
flags compactos; suele combinarse con DP exponencial.}
\complexity{operaciones de bits $O(1)$; enumerar todos los subconjuntos $O(2^n)$
y todas las subm\'ascaras de todas las m\'ascaras $O(3^n)$; memoria $O(1)$.}

Para evitar desplazamientos indefinidos, usar tipos sin signo y verificar que
$0\le k<64$.

\begin{lstlisting}[style=cpp]
using ull = unsigned long long;

// Se debe asegurar que k este en el rango [0,63]
bool has_bit(ull mask, int k) { return (mask >> k) & 1ULL; }
void add_bit(ull &mask, int k) { mask |= 1ULL << k; }
void remove_bit(ull &mask, int k) { mask &= ~(1ULL << k); }
void toggle_bit(ull &mask, int k) { mask ^= 1ULL << k; }
// Las funciones anteriores funcionan con un unico bit
// Sin embargo tambien se puede trabajar con varios bits al mismo tiempo usando lo siguiente
bool has_bits(ull mask, ull bits) { return (ull)(bits & mask) == bits; }
void add_bits(ull &mask, ull bits) { mask |= bits; }
void remove_bits(ull &mask, ull bits) { mask &= ~bits; }
void toggle_bits(ull &mask, ull bits) { mask ^= bits; }

vector<ull> nonempty_submasks(ull mask) {
    vector<ull> answer;
    for (ull sub = mask; sub; sub = (sub - 1) & mask)
        answer.push_back(sub);
    return answer;
}
\end{lstlisting}

Intersecci\'on: \verb+a & b+; uni\'on: \verb+a | b+; diferencia:
\verb+a & ~b+; diferencia sim\'etrica: \verb+a ^ b+.

\section{Conversi\'on a base arbitraria}
\whenuse{Extraer d\'igitos, trabajar con representaciones posicionales o convertir
un entero no negativo a una base entre 2 y 36.}
\complexity{$O(\log_b n)$ tiempo y memoria para producir los d\'igitos.}

Los d\'igitos se devuelven como enteros, por lo que admite bases mayores que
36. El signo se maneja por separado.

\begin{lstlisting}[style=cpp]
pair<bool, vector<int>> to_base(ll value, int base) {
    if (base < 2) throw invalid_argument("base < 2");
    bool negative = value < 0;
    unsigned long long x = negative
        ? (unsigned long long)(-(value + 1)) + 1
        : (unsigned long long)value;

    vector<int> digits;
    do {
        digits.push_back((int)(x % base));
        x /= base;
    } while (x);
    reverse(digits.begin(), digits.end());
    return {negative, digits};
}
\end{lstlisting}
\begin{lstlisting}[style=cpp]
char digit_to_char(int digit) {
    if (digit < 10) return char('0' + digit);
    return char('A' + digit - 10);
}

string to_base_string(ll value, int base) {
    if (base < 2 || base > 36)
        throw invalid_argument("base fuera de [2,36]");

    auto [negative, digits] = to_base(value, base);

    string result;
    if (negative) result += '-';
    for (int digit : digits)
        result += digit_to_char(digit);

    return result;
}
\end{lstlisting}
\begin{lstlisting}[style=cpp]
int char_to_digit(char c) {
    if ('0' <= c && c <= '9') return c - '0';
    if ('A' <= c && c <= 'Z') return c - 'A' + 10;
    if ('a' <= c && c <= 'z') return c - 'a' + 10;
    return -1;
}

ll from_base(const string& text, int base) {
    ll value = 0;
    for (char c : text) {
        int digit = char_to_digit(c);
        if (digit < 0 || digit >= base)
            throw invalid_argument("digito invalido");
        value = value * base + digit;
    }
    return value;
}
\end{lstlisting}